\section{Problemas resueltos de regresión lineal con \texttt{scikit-learn}}

\subsection{Enunciados de los problemas}

\subsubsection*{1. Nivel Fundamental (Sencillos / Directos)}

\begin{problema}
	\label{prob:sklearn_fit_predict_score}
	Un científico de datos ajusta un modelo con \texttt{scikit-learn}:
	\begin{align}
		\texttt{lm = LinearRegression(); lm.fit(X\_train, y\_train)}
	\end{align}
	Explique, en una oración cada uno, qué calcula y qué retorna exactamente cada uno de los siguientes tres métodos: \texttt{lm.fit(X,y)}, \texttt{lm.predict(X\_new)}, y \texttt{lm.score(X,y)}. Especifique cuál de los tres modifica el estado interno del objeto \texttt{lm} y cuáles no.
\end{problema}

\begin{problema}
	\label{prob:sklearn_coef_intercept}
	Al ajustar un modelo \texttt{LinearRegression} con dos predictores (\texttt{TV}, \texttt{Radio}) sobre datos de ventas, se obtiene:
	\begin{align}
		\texttt{lm.intercept\_} = 2.92, \qquad \texttt{lm.coef\_} = [0.046,\ 0.188].
	\end{align}
	Escriba la ecuación del modelo en notación algebraica estándar y calcule la venta predicha para un presupuesto de \texttt{TV}$=100$ y \texttt{Radio}$=20$.
\end{problema}

\begin{problema}
	\label{prob:sklearn_train_test_split_basico}
	Un analista tiene $n=250$ observaciones y ejecuta \texttt{train\_test\_split(X, y, test\_size=0.2, random\_state=42)}. Calcule el número exacto de observaciones en los conjuntos de entrenamiento y de prueba, y explique por qué fijar \texttt{random\_state} garantiza la reproducibilidad exacta de la partición entre distintas ejecuciones del código.
\end{problema}

\subsubsection*{2. Nivel Operativo (Intermedios / De Desarrollo)}

\begin{problema}
	\label{prob:sklearn_r2_manual_vs_score}
	Sobre un conjunto de prueba de $n=5$ observaciones, un modelo predice $\hat Y=\{18.2,\ 9.5,\ 8.1,\ 14.0,\ 15.9\}$ contra los valores reales $Y=\{17.9,\ 10.4,\ 7.8,\ 14.2,\ 15.6\}$. Calcule manualmente el $R^2$ usando $\text{SCR}=\sum(\hat Y_i-\bar Y)^2$ y $\text{SCT}=\sum(Y_i-\bar Y)^2$, y explique por qué este valor debe coincidir exactamente con lo que retornaría \texttt{lm.score(X\_test, y\_test)}.
\end{problema}

\begin{problema}
	\label{prob:sklearn_rfe_interpretacion}
	Al ejecutar \texttt{RFE(estimator, n\_features\_to\_select=2)} sobre 4 predictores (\texttt{TV}, \texttt{Radio}, \texttt{Newspaper}, \texttt{Online}), se obtiene:
	\begin{align}
		\texttt{selector.support\_} = [\texttt{True},\texttt{True},\texttt{False},\texttt{False}], \qquad \texttt{selector.ranking\_} = [1,1,3,2].
	\end{align}
	Determine qué predictores fueron seleccionados para el modelo final, y determine el orden exacto en que los predictores no seleccionados fueron eliminados durante el proceso recursivo.
\end{problema}

\begin{problema}
	\label{prob:sklearn_fit_intercept_false}
	Un analista ajusta accidentalmente \texttt{LinearRegression(fit\_intercept=False)} sobre un conjunto de datos donde la variable respuesta \texttt{Sales} nunca es cero cuando los predictores son cero (existe un piso de ventas orgánicas). Explique matemáticamente qué restricción geométrica impone \texttt{fit\_intercept=False} sobre la recta ajustada, y por qué esto probablemente sesgará las estimaciones de las pendientes $\hat\beta_j$ en este escenario.
\end{problema}

\subsubsection*{3. Nivel Analítico (Avanzados / De Conexión)}

\begin{problema}
	\label{prob:sklearn_svd_vs_ecuacion_normal}
	\texttt{scikit-learn} no resuelve \texttt{LinearRegression} invirtiendo directamente $\mathbf{X}^T\mathbf{X}$ como en la Sección 09.07; en su lugar, utiliza una descomposición en valores singulares (SVD) de $\mathbf{X}$. Explique por qué ambos métodos producen el mismo resultado cuando $\mathbf{X}$ tiene rango columna completo, y explique la ventaja numérica de la SVD cuando $\mathbf{X}^T\mathbf{X}$ está cerca de ser singular por multicolinealidad severa (Sección 09.09).
\end{problema}

\begin{problema}
	\label{prob:sklearn_rfe_costo_computacional}
	Considere \texttt{RFE(estimator, n\_features\_to\_select=k, step=1)} aplicado sobre un conjunto de $p$ predictores totales ($k<p$). En cada iteración, el algoritmo ajusta el estimador sobre el subconjunto de predictores actualmente activos y elimina exactamente 1 predictor (el de menor importancia) antes de continuar. Deduzca una fórmula general para el número total de modelos ajustados internamente por \texttt{RFE} durante todo el proceso, en función de $p$ y $k$.
\end{problema}

\subsubsection*{4. Nivel Desafiante (Expertos / Deducción Profunda)}

\begin{problema}
	\label{prob:sklearn_pseudoinversa_deduccion}
	\textbf{Deducción de la equivalencia entre la solución vía SVD y la Ecuación Normal:}
	Sea $\mathbf{X}=\mathbf{U}\boldsymbol\Sigma\mathbf{V}^T$ la descomposición en valores singulares de la matriz de diseño ($\mathbf{U}$ y $\mathbf{V}$ ortogonales, $\boldsymbol\Sigma$ diagonal con las singulares $\sigma_i>0$). Demuestre que la solución de mínimos cuadrados obtenida mediante la pseudoinversa de Moore-Penrose $\hat{\boldsymbol\beta}=\mathbf{X}^+\mathbf{Y}$, donde $\mathbf{X}^+=\mathbf{V}\boldsymbol\Sigma^{-1}\mathbf{U}^T$, es algebraicamente idéntica a $\hat{\boldsymbol\beta}=(\mathbf{X}^T\mathbf{X})^{-1}\mathbf{X}^T\mathbf{Y}$ cuando $\mathbf{X}$ tiene rango columna completo.
\end{problema}

\begin{problema}
	\label{prob:sklearn_multicolinealidad_severa_svd}
	\textbf{Análisis crítico de la solución de mínima norma bajo colinealidad perfecta:}
	Suponga que dos predictores del modelo son perfectamente colineales ($X_2=2X_1$ exactamente), de modo que $\mathbf{X}^T\mathbf{X}$ es exactamente singular (no invertible) y la Ecuación Normal de la Sección 09.07 no tiene solución única. Explique por qué \texttt{scikit-learn} (usando SVD internamente) aun así retorna un vector de coeficientes $\hat{\boldsymbol\beta}$ sin error, qué representa exactamente esa solución particular (la de \emph{norma mínima} entre todas las soluciones que minimizan $\|\mathbf{Y}-\mathbf{X}\boldsymbol\beta\|^2$), y por qué esta solución, aunque numéricamente válida, no debe interpretarse como si cada coeficiente individual $\hat\beta_j$ tuviera un significado causal aislado.
\end{problema}


\subsection{Soluciones detalladas}

\subsubsection*{1. Nivel Fundamental (Sencillos / Directos)}

\begin{solproblema}
	\texttt{lm.fit(X,y)} calcula los coeficientes óptimos de mínimos cuadrados y los almacena internamente en \texttt{lm.intercept\_} y \texttt{lm.coef\_}, modificando el estado del objeto; \texttt{lm.predict(X\_new)} usa esos coeficientes ya almacenados para calcular $\hat Y=\hat\beta_0+\hat\beta_1X_{\text{new}}$ sin modificar el estado del objeto; \texttt{lm.score(X,y)} calcula el $R^2$ del modelo ya ajustado sobre el conjunto $(X,y)$ proporcionado, tampoco modifica el estado. Solo \texttt{fit} altera el objeto \texttt{lm}.
\end{solproblema}

\begin{solproblema}
	La ecuación del modelo es $\hat{\texttt{Sales}}=2.92+0.046\,\texttt{TV}+0.188\,\texttt{Radio}$. Sustituyendo \texttt{TV}$=100$, \texttt{Radio}$=20$:
	\begin{align}
		\hat{\texttt{Sales}}=2.92+0.046(100)+0.188(20)=2.92+4.6+3.76=11.28.
	\end{align}
\end{solproblema}

\begin{solproblema}
	Con $n=250$ y $\texttt{test\_size}=0.2$: $n_{\text{test}}=0.2\times250=50$, $n_{\text{train}}=250-50=200$. Fijar \texttt{random\_state} inicializa el generador de números pseudoaleatorios con una semilla determinista: el algoritmo de mezcla y partición interno se vuelve completamente reproducible, generando exactamente la misma asignación de observaciones a cada subconjunto cada vez que se ejecuta el código con la misma semilla.
\end{solproblema}

\subsubsection*{2. Nivel Operativo (Intermedios / De Desarrollo)}

\begin{solproblema}
	Calculamos $\bar Y=(17.9+10.4+7.8+14.2+15.6)/5=13.18$.
	\begin{align}
		\text{SCT}=\sum(Y_i-\bar Y)^2=(4.72)^2+(2.78)^2+(5.38)^2+(1.02)^2+(2.42)^2\approx22.28+7.73+28.94+1.04+5.86\approx65.85
	\end{align}
	\begin{align}
		\text{SCR}=\sum(\hat Y_i-\bar Y)^2=(5.02)^2+(3.68)^2+(5.08)^2+(0.82)^2+(2.72)^2\approx25.20+13.54+25.81+0.67+7.40\approx72.62
	\end{align}
	Dado que aquí se solicita $R^2=1-\text{SCE}/\text{SCT}$ con $\text{SCE}=\sum(Y_i-\hat Y_i)^2=(0.3)^2+(-0.9)^2+(0.3)^2+(0.2)^2+(0.3)^2=0.09+0.81+0.09+0.04+0.09=1.12$:
	\begin{align}
		R^2=1-\frac{1.12}{65.85}\approx0.983.
	\end{align}
	Este valor debe coincidir exactamente con \texttt{lm.score(X\_test,y\_test)} porque \texttt{score()} implementa precisamente esta misma fórmula ($R^2=1-\text{SCE}/\text{SCT}$) internamente; no existe ningún cálculo alternativo oculto en el método.
\end{solproblema}

\begin{solproblema}
	Los predictores seleccionados son aquellos con \texttt{support\_}$=\texttt{True}$: \texttt{TV} y \texttt{Radio}. Para los eliminados, un ranking menor indica eliminación más tardía (mayor importancia relativa); ranking mayor indica eliminación más temprana. Con \texttt{ranking\_}$=[1,1,3,2]$: \texttt{Online} (ranking 2) fue eliminado en la penúltima ronda, y \texttt{Newspaper} (ranking 3) fue el primero en ser eliminado, siendo el predictor considerado menos relevante desde el inicio del proceso recursivo.
\end{solproblema}

\begin{solproblema}
	Geométricamente, \texttt{fit\_intercept=False} obliga a la recta o hiperplano ajustado a pasar exactamente por el origen ($\hat Y=0$ cuando todos los predictores son $0$). Si en la realidad existe un nivel base de \texttt{Sales} positivo incluso sin inversión publicitaria, forzar el modelo a pasar por el origen introduce un sesgo sistemático: el algoritmo de mínimos cuadrados compensará esa restricción incorrecta inflando artificialmente las pendientes $\hat\beta_j$ para intentar explicar ese nivel base a través de los predictores, en lugar de a través de un intercepto legítimamente distinto de cero.
\end{solproblema}

\subsubsection*{3. Nivel Analítico (Avanzados / De Conexión)}

\begin{solproblema}
	Cuando $\mathbf{X}$ tiene rango columna completo, la matriz $\mathbf{X}^T\mathbf{X}$ es invertible y ambos métodos --- inversión directa y SVD --- resuelven exactamente el mismo sistema de ecuaciones normales, produciendo resultados idénticos hasta el error de redondeo de punto flotante. La ventaja de la SVD surge cuando $\mathbf{X}^T\mathbf{X}$ está cerca de ser singular: el número de condición de $\mathbf{X}^T\mathbf{X}$ es aproximadamente el \textbf{cuadrado} del número de condición de $\mathbf{X}$, de modo que invertir $\mathbf{X}^T\mathbf{X}$ directamente amplifica dramáticamente los errores numéricos de redondeo. La SVD opera directamente sobre $\mathbf{X}$ (no sobre $\mathbf{X}^T\mathbf{X}$), evitando ese cuadrado del número de condición y produciendo una solución numéricamente mucho más estable ante multicolinealidad severa.
\end{solproblema}

\begin{solproblema}
	Partiendo de $p$ predictores y terminando en $k$, con eliminación de exactamente 1 predictor por iteración, el número de iteraciones de eliminación es $p-k$. En cada una de esas iteraciones se ajusta un modelo sobre el subconjunto de predictores actualmente activo, y se requiere un ajuste final adicional sobre el subconjunto final de $k$ predictores seleccionados. Por lo tanto, el número total de modelos ajustados es:
	\begin{align}
		N_{\text{ajustes}} = (p-k)+1.
	\end{align}
	Por ejemplo, para $p=10$ predictores reduciendo a $k=3$, \texttt{RFE} ajusta $10-3+1=8$ modelos en total.
\end{solproblema}

\subsubsection*{4. Nivel Desafiante (Expertos / Deducción Profunda)}

\begin{solproblema}
	\textbf{Demostración de la equivalencia SVD--Ecuación Normal:}

	Con $\mathbf{X}=\mathbf{U}\boldsymbol\Sigma\mathbf{V}^T$, calculamos $\mathbf{X}^T\mathbf{X}$:
	\begin{align}
		\mathbf{X}^T\mathbf{X} = (\mathbf{U}\boldsymbol\Sigma\mathbf{V}^T)^T(\mathbf{U}\boldsymbol\Sigma\mathbf{V}^T) = \mathbf{V}\boldsymbol\Sigma\mathbf{U}^T\mathbf{U}\boldsymbol\Sigma\mathbf{V}^T = \mathbf{V}\boldsymbol\Sigma^2\mathbf{V}^T,
	\end{align}
	usando $\mathbf{U}^T\mathbf{U}=\mathbf{I}$ (ortogonalidad de $\mathbf{U}$). Como $\mathbf{V}$ es ortogonal ($\mathbf{V}^{-1}=\mathbf{V}^T$) y $\boldsymbol\Sigma^2$ es diagonal invertible (rango completo, todos los $\sigma_i>0$), la inversa es:
	\begin{align}
		(\mathbf{X}^T\mathbf{X})^{-1} = \mathbf{V}\boldsymbol\Sigma^{-2}\mathbf{V}^T.
	\end{align}
	Ahora calculamos $\mathbf{X}^T\mathbf{Y} = (\mathbf{U}\boldsymbol\Sigma\mathbf{V}^T)^T\mathbf{Y} = \mathbf{V}\boldsymbol\Sigma\mathbf{U}^T\mathbf{Y}$. Sustituyendo en la Ecuación Normal:
	\begin{align}
		\hat{\boldsymbol\beta} = (\mathbf{X}^T\mathbf{X})^{-1}\mathbf{X}^T\mathbf{Y} = \mathbf{V}\boldsymbol\Sigma^{-2}\mathbf{V}^T\mathbf{V}\boldsymbol\Sigma\mathbf{U}^T\mathbf{Y} = \mathbf{V}\boldsymbol\Sigma^{-2}\boldsymbol\Sigma\mathbf{U}^T\mathbf{Y} = \mathbf{V}\boldsymbol\Sigma^{-1}\mathbf{U}^T\mathbf{Y},
	\end{align}
	usando $\mathbf{V}^T\mathbf{V}=\mathbf{I}$. Como $\mathbf{X}^+=\mathbf{V}\boldsymbol\Sigma^{-1}\mathbf{U}^T$ por definición de la pseudoinversa de Moore-Penrose, concluimos exactamente:
	\begin{align}
		\hat{\boldsymbol\beta} = \mathbf{X}^+\mathbf{Y} = (\mathbf{X}^T\mathbf{X})^{-1}\mathbf{X}^T\mathbf{Y}.
	\end{align}
	Las dos expresiones son algebraicamente idénticas cuando $\mathbf{X}$ tiene rango columna completo (todos los $\sigma_i>0$, garantizando que $\boldsymbol\Sigma^{-1}$ existe).
\end{solproblema}

\begin{solproblema}
	Cuando $X_2=2X_1$ exactamente, las columnas de $\mathbf{X}$ son linealmente dependientes: $\mathbf{X}$ no tiene rango columna completo, al menos un valor singular $\sigma_i=0$, y $\mathbf{X}^T\mathbf{X}$ es exactamente singular (determinante cero, sin inversa). La Ecuación Normal de la Sección 09.07 no tiene solución única: existen infinitas combinaciones de $(\hat\beta_1,\hat\beta_2)$ que producen la misma predicción $\hat Y$, ya que cualquier incremento en $\hat\beta_1$ compensado por un decremento apropiado en $\hat\beta_2$ (manteniendo $\hat\beta_1+2\hat\beta_2$ constante) deja $\mathbf{X}\hat{\boldsymbol\beta}$ inalterado.

	\texttt{scikit-learn} no falla porque la pseudoinversa de Moore-Penrose $\mathbf{X}^+$ está definida incluso cuando algún $\sigma_i=0$ (simplemente se omite ese término, o se usa $\sigma_i^{-1}=0$ por convención en la descomposición truncada). La solución retornada, $\hat{\boldsymbol\beta}=\mathbf{X}^+\mathbf{Y}$, es específicamente aquella de \textbf{norma euclidiana mínima} ($\|\hat{\boldsymbol\beta}\|_2$ mínima) entre todas las soluciones que logran el mismo ajuste óptimo $\|\mathbf{Y}-\mathbf{X}\hat{\boldsymbol\beta}\|^2$ mínimo.

	Esta solución es numéricamente válida y reproducible, pero \textbf{no debe interpretarse causalmente coeficiente por coeficiente}: dado que $(\hat\beta_1,\hat\beta_2)$ es solo una elección arbitraria (la de norma mínima) entre infinitas combinaciones equivalentes, el valor específico de $\hat\beta_1$ no representa de forma aislada e inequívoca "el efecto de $X_1$ manteniendo $X_2$ constante" --- esa interpretación pierde sentido precisamente porque $X_1$ y $X_2$ nunca varían de forma independiente en los datos observados.
\end{solproblema}
